\section{Análisis}

Prueba corrida con \texttt{visor\_enlace} (interfaz gráfica en la misma PC),
con el mismo firmware que \texttt{resultados/2026-10-05\_pll}
($F_{cy} = 50$\,MHz). El patrón de errores es el mismo que con
\texttt{prueba\_enlace}: no hay errores de encabezado ni reinicios, y cada
respuesta con CRC incorrecto va seguida de la bandera de longitud del dsPIC
para la misma transferencia (columna ``status siguiente'' en la captura). La
corrupción empieza en el byte 1 en el 70\,\% de los casos y en el byte 7 en
el 20\,\% (61\,\% y 26\,\% con \texttt{prueba\_enlace}).

La proporción de tramas afectadas fue de 1.0\,\%, contra 0.7\,\% con
\texttt{prueba\_enlace}. Con dos corridas no se puede distinguir si la
diferencia se debe a la carga de la interfaz gráfica o a la variación entre
corridas; el origen del defecto está en la frontera de bytes del SPI esclavo
y debe revisarse con el analizador lógico.
